%==============================================================================
% Chapitre 14 - Introduction aux armes à feu
%==============================================================================

\section{Introduction}

Une arme à feu est un système mécanique destiné à propulser un projectile
à grande vitesse à l'aide de l'énergie libérée par la combustion d'une
charge propulsive.

Depuis plusieurs siècles, les armes à feu constituent l'un des principaux
moyens de projection de projectiles à moyenne et longue distance.

Elles sont utilisées dans des domaines variés :

\begin{itemize}
    \item militaire ;
    \item sécurité intérieure ;
    \item maintien de l'ordre ;
    \item chasse ;
    \item tir sportif ;
    \item essais et recherche.
\end{itemize}

\begin{aretenir}

Une arme à feu transforme l'énergie chimique contenue dans une munition en
énergie cinétique transmise à un projectile.

\end{aretenir}

\section{Vue d'ensemble}

Quelle que soit sa complexité, une arme à feu remplit toujours les mêmes
fonctions fondamentales :

\begin{enumerate}
    \item recevoir une munition ;
    \item initier la combustion de la charge propulsive ;
    \item guider le projectile ;
    \item évacuer les gaz ;
    \item permettre le contrôle du tir.
\end{enumerate}

Les solutions techniques utilisées pour réaliser ces fonctions peuvent être
très différentes selon le type d'arme considéré.

\section{Principe général de fonctionnement}

Lors du tir :

\begin{enumerate}
    \item une munition est introduite dans la chambre ;
    \item l'amorce est percutée ;
    \item la charge propulsive s'enflamme ;
    \item la pression augmente rapidement ;
    \item le projectile est accéléré dans le canon ;
    \item le projectile quitte l'arme.
\end{enumerate}

Le cycle complet sera étudié en détail dans la partie consacrée à la
balistique intérieure.

\section{Transformation d'énergie}

Une arme à feu constitue essentiellement un convertisseur d'énergie.

L'énergie suit généralement la chaîne suivante :

\begin{center}
Énergie chimique
$\rightarrow$
Énergie thermique
$\rightarrow$
Énergie mécanique
$\rightarrow$
Énergie cinétique
\end{center}

Une partie de l'énergie disponible est toutefois perdue sous forme :

\begin{itemize}
    \item de chaleur ;
    \item de bruit ;
    \item de vibrations ;
    \item de lumière ;
    \item de mouvements parasites.
\end{itemize}

\section{Les principaux sous-systèmes}

Une arme moderne peut être décomposée en plusieurs sous-ensembles.

Les plus importants sont :

\begin{itemize}
    \item le canon ;
    \item la chambre ;
    \item le mécanisme de verrouillage ;
    \item le système d'alimentation ;
    \item le mécanisme de détente ;
    \item les organes de visée ;
    \item la structure porteuse.
\end{itemize}

Chacun de ces éléments possède une influence directe sur les performances
du système.

\section{Classification des armes}

Les armes à feu peuvent être classées selon différents critères.

\subsection{Selon le mode de fonctionnement}

On distingue notamment :

\begin{itemize}
    \item les armes à répétition manuelle ;
    \item les armes semi-automatiques ;
    \item les armes automatiques.
\end{itemize}

\subsection{Selon le type d'emploi}

On rencontre par exemple :

\begin{itemize}
    \item les armes de poing ;
    \item les fusils ;
    \item les fusils de précision ;
    \item les mitrailleuses ;
    \item les canons automatiques.
\end{itemize}

\subsection{Selon le calibre}

Le calibre constitue l'une des caractéristiques les plus utilisées pour
décrire une arme.

Il sera étudié plus en détail dans les chapitres suivants.

\section{Performances d'une arme}

Les performances d'une arme ne se limitent pas à la vitesse du projectile.

D'autres paramètres sont également importants :

\begin{itemize}
    \item précision ;
    \item dispersion ;
    \item cadence de tir ;
    \item portée ;
    \item fiabilité ;
    \item ergonomie ;
    \item masse.
\end{itemize}

Le choix d'un système d'arme résulte souvent d'un compromis entre ces
différents critères.

\section{Armes et balistique}

L'arme influence directement plusieurs domaines de la balistique.

Elle intervient notamment dans :

\begin{itemize}
    \item la vitesse initiale ;
    \item la stabilité du projectile ;
    \item la précision ;
    \item le recul ;
    \item les vibrations ;
    \item la dispersion.
\end{itemize}

Les caractéristiques du canon et du système de tir jouent un rôle majeur
dans les performances finales.

\section{Armes réelles et simulation}

Dans une simulation, une arme est généralement représentée par un ensemble
de paramètres.

Selon le niveau de fidélité recherché, ces paramètres peuvent inclure :

\begin{itemize}
    \item la longueur du canon ;
    \item le calibre ;
    \item la vitesse initiale ;
    \item la dispersion ;
    \item la cadence de tir ;
    \item les caractéristiques des munitions.
\end{itemize}

Les modèles les plus avancés prennent également en compte :

\begin{itemize}
    \item les vibrations ;
    \item l'échauffement ;
    \item l'usure ;
    \item les mouvements de la plateforme.
\end{itemize}

\begin{simulation}

Dans de nombreux simulateurs, l'arme est représentée de manière simplifiée.

Les modèles les plus réalistes cherchent à reproduire le comportement du
système complet, depuis le départ du coup jusqu'à l'impact du projectile.

\end{simulation}

\section{Vers les chapitres suivants}

Les prochains chapitres détailleront les principaux composants des armes à
feu.

Nous étudierons notamment :

\begin{itemize}
    \item le canon ;
    \item la chambre ;
    \item les mécanismes ;
    \item les systèmes d'alimentation ;
    \item les dispositifs de visée.
\end{itemize}

Cette étude permettra ensuite de comprendre l'origine de nombreux phénomènes
balistiques observés lors du tir.

\section{À retenir}

\begin{aretenir}

\begin{itemize}
    \item Une arme à feu transforme une énergie chimique en énergie
    cinétique.
    \item Le canon guide et accélère le projectile.
    \item Une arme moderne est constituée de plusieurs sous-systèmes.
    \item Les caractéristiques de l'arme influencent directement les
    performances balistiques.
    \item L'étude détaillée des composants est nécessaire pour comprendre le
    comportement global du système.
\end{itemize}

\end{aretenir}
